\chapter{晶体的定义}
\markboth{第 2 章\ 晶体的定义}{第 2 章\ 晶体的定义}

本节定义左晶体（对 $\PreStk$ 的任意对象），以及对 $\PreStk_{\mathrm{laft}}$ 的对象定义右晶体。我们将证明：在后一情形，两套理论等价。

\section{左晶体}

\ob{2.1.1}{n:2.1.1}对 $Y\in\PreStk$，定义
\[
\Crys^l(Y):=\QCoh(Y_{\dR}).
\]
也就是说，
\[
\Crys^l(Y)=\varprojlim_{S\in(\DGSch_{\mathrm{aff}}/Y_{\dR})^{\mathrm{op}}}\QCoh(S).
\]
粗略地说，对象 $M\in\Crys^l(Y)$ 是如下资料：对每个带映射 $\clred S\to Y$ 的仿射 DG 概形 $S\in\DGSch_{\mathrm{aff}}$，指定一个拟凝聚层 $F_S\in\QCoh(S)$；并且对每个仿射 DG 概形的态射 $f:S'\to S$，指定同构
\[
f^*(F_S)\simeq F_{S'}\in\QCoh(S').
\]

\ob{2.1.2}{n:2.1.2}更函子化地说，记 $\Crys^l_{\PreStk}$ 为函子 $(\PreStk)^{\mathrm{op}}\to\DGCat_{\mathrm{cont}}$，定义为
\[
\Crys^l_{\PreStk}:=\QCoh^*_{\PreStk}\circ\dR,
\]
其中
\[
\QCoh^*_{\PreStk}:(\PreStk)^{\mathrm{op}}\to\DGCat_{\mathrm{cont}}
\]
是把预栈对应到其拟凝聚层范畴的函子 [GL:QCoh, Sect.~1.1.5]。

对 $\PreStk$ 中的映射 $f:Y_1\to Y_2$，记 $f^{\dagger,l}$ 为相应的拉回函子
\[
\Crys^l(Y_2)\to\Crys^l(Y_1).
\]
按构造，若 $f$ 诱导底层既约经典预栈的同构 $\clred Y_1\to\clred Y_2$，则它诱导 de Rham 预栈的同构 $Y_{1,\dR}\to Y_{2,\dR}$；特别地，$f^{\dagger,l}$ 是等价。

\ob{2.1.3}{n:2.1.3}回忆函子 $\QCoh^*_{\PreStk}:(\PreStk)^{\mathrm{op}}\to\DGCat_{\mathrm{cont}}$ 按定义是函子
\[
\QCoh^*_{\DGSch_{\mathrm{aff}}}:(\DGSch_{\mathrm{aff}})^{\mathrm{op}}\to\DGCat_{\mathrm{cont}}
\]
沿 $(\DGSch_{\mathrm{aff}})^{\mathrm{op}}\hookrightarrow(\PreStk)^{\mathrm{op}}$ 的右 Kan 延拓。

特别地，它把 $\PreStk$ 中的余极限送到 $\DGCat_{\mathrm{cont}}$ 中的极限。因此由推论 1.1.9，对 $Y\in\PreStk$ 得到：

\begin{othm}{推论}{2.1.4}{n:2.1.4}
设 $\mathcal{C}$ 是 §1.1.8 列表中的任一范畴。则对 $Y\in\PreStk$，函子
\[
\Crys^l(Y)\to\varprojlim_{X\in(\mathcal{C}_{/Y})^{\mathrm{op}}}\Crys^l(X)
\]
是等价。
\end{othm}

粗略地说，这个推论表明：对象 $M\in\Crys^l(Y)$ 的数据，等价于对每个 $S\in\mathcal{C}_{/Y}$ 给定 $M_S\in\Crys^l(S)$，并对每个 $f:S'\to S$ 给定同构
\[
f^{\dagger,l}(M_S)\simeq M_{S'}\in\Crys^l(S').
\]

\ob{2.1.5}{n:2.1.5}回忆自然变换 $p_{\dR}:\mathrm{Id}\to\dR$。它诱导自然变换
\[
\oblv^l:\Crys^l_{\PreStk}\to\QCoh^*_{\PreStk}.
\]
也就是说，对每个 $Y\in\PreStk$ 有函子
\begin{equation}
\oblv^l_Y:\Crys^l(Y)\to\QCoh(Y),\tag{2.1}
\end{equation}
且对每个态射 $f:Y_1\to Y_2$ 有交换图
\begin{equation}
\begin{tikzcd}[column sep=large, row sep=large]
 \Crys^l(Y_1) \arrow[r, "\oblv^l_{Y_1}"] \arrow[d, "f^{\dagger,l}"'] & \QCoh(Y_1) \arrow[d, "f^*"] \\
 \Crys^l(Y_2) \arrow[r, "\oblv^l_{Y_2}"'] & \QCoh(Y_2).
\end{tikzcd}\tag{2.2}
\end{equation}

\ob{2.1.6}{n:2.1.6}回忆 §1.2.2 的单纯对象 $Y_\bullet/Y_{\dR}$。

由引理 1.2.4 得到：

\begin{othm}{引理}{2.1.7}{n:2.1.7}
若 $Y$ 经典形式光滑，则函子
\[
\Crys^l(Y)\to\Tot\big(\QCoh(Y_\bullet/Y_{\dR})\big)
\]
是等价。
\end{othm}

\begin{othm}{注}{2.1.8}{n:2.1.8}
本书对 $Y$ 上左晶体的定义，就是 Grothendieck 术语中 $Y$ 的无穷小 site 上的拟凝聚层。范畴 $\Tot(\QCoh(Y_\bullet/Y_{\dR}))$ 则是 Grothendieck 术语中 $Y$ 的分层 site 上的拟凝聚层。于是引理 2.1.7 说：对经典形式光滑的预栈，两者等价。我们将在 §3.4 看到，当 $Y$ 是局部几乎有限型的最终余连接 DG 概形时，同样的结论也成立。不过，对非最终余连接的 DG 概形（即使是局部几乎有限型的），这个等价不成立。
\end{othm}

\section{局部几乎有限型预栈上的左晶体}

在本节余下部分，除另有说明，我们把讨论限制在子范畴 $\PreStk_{\mathrm{laft}}\subset\PreStk$。

因此，除明确说明外，所谓预栈/DG 概形/仿射 DG 概形，都指局部几乎有限型的。

记 $\Crys^l_{\PreStk_{\mathrm{laft}}}$ 为 $\Crys^l_{\PreStk}$ 到 $\PreStk_{\mathrm{laft}}\subset\PreStk$ 的限制。

\ob{2.2.1}{n:2.2.1}下面的推论说明：为了对 $Y\in\PreStk_{\mathrm{laft}}$ 定义 $\Crys^l(Y)$，我们「不需要知道」无限型概形或导出代数几何。换言之，要在局部几乎有限型的预栈上定义晶体，可以留在有限型经典仿射概形的世界里。

事实上，由推论 1.3.7 得到：

\begin{othm}{推论}{2.2.2}{n:2.2.2}
设 $\mathcal{C}$ 是 $\DGSch_{\mathrm{aff}}$ 的下列全子范畴之一：
\[
\Sch_{\mathrm{aff}}^{\mathrm{ft}},\quad {}^{<\infty}\DGSch_{\mathrm{aff}}^{\mathrm{aft}},\quad \DGSch_{\mathrm{aff}}^{\mathrm{aft}},\quad \Sch_{\mathrm{aff}}.
\]
则对 $Y\in\PreStk_{\mathrm{laft}}$，自然函子
\[
\Crys^l(Y)\to\varprojlim_{S\in(\mathcal{C}_{/Y_{\dR}})^{\mathrm{op}}}\QCoh(S)
\]
是等价。
\end{othm}

\ob{2.2.3}{n:2.2.3}回忆：按推论 2.1.4，范畴 $\Crys^l(Y)$ 可以由函子
\[
\Crys^l:\mathcal{C}_{/Y}\to\DGCat_{\mathrm{cont}}
\]
还原，其中 $\mathcal{C}$ 是下列范畴之一：
\[
{}^{\mathrm{red}}\Sch_{\mathrm{aff}},\ \Sch_{\mathrm{aff}},\ \DGSch_{\mathrm{aff}},\ \DGSch_{\mathrm{qs\text{-}qc}},\ \DGSch\subset\PreStk.
\]
现在断言：上述范畴也可以换成它们在 (1.3) 中的全子范畴：
\[
{}^{\mathrm{red}}\Sch_{\mathrm{aff}}^{\mathrm{ft}},\quad \Sch_{\mathrm{aff}}^{\mathrm{ft}},\quad \DGSch_{\mathrm{aff}}^{\mathrm{ft}},\quad \DGSch_{\mathrm{aff}}^{\mathrm{aft}},\quad \DGSch_{\mathrm{laft}},\quad \PreStk_{\mathrm{laft}}.
\]

\begin{othm}{推论}{2.2.4}{n:2.2.4}
设 $Y\in\PreStk_{\mathrm{laft}}$，$\mathcal{C}$ 是 (1.3) 中的范畴之一。则函子
\[
\Crys^l(Y)\to\varprojlim_{X\in(\mathcal{C}_{/Y})^{\mathrm{op}}}\Crys^l(X)
\]
是等价。
\end{othm}

\begin{proof}
由推论 1.3.9 得到。
\end{proof}

粗略地说，上述推论表明：对象 $M\in\Crys^l(Y)$ 可以由如下资料还原：对每个 $S\in\mathcal{C}_{/Y}$ 给定 $M_S\in\Crys^l(S)$，并对每个 $f:S'\to S$ 给定同构
\[
f^{\dagger,l}(M_S)\simeq M_{S'}\in\Crys^l(S').
\]

\ob{2.2.5}{n:2.2.5}再次考虑 (2.1) 中的函子
\[
\oblv^l_Y:\Crys^l(Y)\to\QCoh(Y).
\]
有：

\begin{othm}{引理}{2.2.6}{n:2.2.6}
对 $Y\in\PreStk_{\mathrm{laft}}$，函子 $\oblv^l_Y$ 保守。
\end{othm}

其证明推迟到 §2.4.7。

\section{右晶体}

\ob{2.3.1}{n:2.3.1}回忆 $\PreStk_{\mathrm{laft}}$ 也可以看作全体函子
\[
({}^{<\infty}\DGSch_{\mathrm{aff}}^{\mathrm{aft}})^{\mathrm{op}}\to\infty\text{-}\Grpd
\]
构成的范畴，见 [GL:Stacks, Sect.~1.3.11]。

此外，有 [IndCoh, Sect.~10.1.2] 中的函子
\[
\IndCoh^!_{\PreStk_{\mathrm{laft}}}:(\PreStk_{\mathrm{laft}})^{\mathrm{op}}\to\DGCat_{\mathrm{cont}},
\]
它定义为相应函子
\[
\IndCoh^!_{{}^{<\infty}\DGSch_{\mathrm{aff}}^{\mathrm{aft}}}:({}^{<\infty}\DGSch_{\mathrm{aff}}^{\mathrm{aft}})^{\mathrm{op}}\to\DGCat_{\mathrm{cont}}
\]
沿
\[
({}^{<\infty}\DGSch_{\mathrm{aff}}^{\mathrm{aft}})^{\mathrm{op}}\to(\PreStk_{\mathrm{laft}})^{\mathrm{op}}
\]
的右 Kan 延拓。

特别地，函子 $\IndCoh^!_{\PreStk_{\mathrm{laft}}}$ 把 $\PreStk_{\mathrm{laft}}$ 中的余极限送到 $\DGCat_{\mathrm{cont}}$ 中的极限。

\ob{2.3.2}{n:2.3.2}定义函子
\[
\Crys^r_{\PreStk_{\mathrm{laft}}}:(\PreStk_{\mathrm{laft}})^{\mathrm{op}}\to\DGCat_{\mathrm{cont}}
\]
为复合
\[
\Crys^r_{\PreStk_{\mathrm{laft}}}:=\IndCoh^!_{\PreStk_{\mathrm{laft}}}\circ\dR.
\]
上式中用到命题 1.3.3(a)，以确保 $\dR$ 是函子 $\PreStk_{\mathrm{laft}}\to\PreStk_{\mathrm{laft}}$。

\begin{othm}{注}{2.3.3}{n:2.3.3}
在定义 $\Crys^r_{\PreStk_{\mathrm{laft}}}$ 时，我们「不需要知道」无限型概形：自函子 $\dR:\PreStk_{\mathrm{laft}}\to\PreStk_{\mathrm{laft}}$ 可以直接由公式
\[
\Maps(S,Y_{\dR})=\Maps(\clred S,Y)
\]
定义，其中 $S\in{}^{<\infty}\DGSch_{\mathrm{aff}}^{\mathrm{aft}}$。
\end{othm}

\ob{2.3.4}{n:2.3.4}对 $\PreStk_{\mathrm{laft}}$ 中的映射 $f:Y_1\to Y_2$，记 $f^{\dagger,r}$ 为相应的函子 $\Crys^r(Y_2)\to\Crys^r(Y_1)$。

若 $f$ 诱导等价 $\clred Y_1\to\clred Y_2$，则映射 $Y_{1,\dR}\to Y_{2,\dR}$ 是等价；特别地，$f^{\dagger,r}$ 也是等价。

\ob{2.3.5}{n:2.3.5}按定义，对 $Y\in\PreStk_{\mathrm{laft}}$ 有
\[
\Crys^r(Y)=\varprojlim_{S\in(({}^{<\infty}\DGSch_{\mathrm{aff}}^{\mathrm{aft}})_{/Y_{\dR}})^{\mathrm{op}}}\IndCoh(S).
\]
粗略地说，对象 $M\in\Crys^r(Y)$ 是如下资料：对每个 $S\in{}^{<\infty}\DGSch_{\mathrm{aff}}^{\mathrm{aft}}$ 及映射 $\clred S\to Y$，指定一个对象 $F_S\in\IndCoh(S)$；并且对每个 $f:S'\to S$，指定同构
\[
f^!(F_S)\simeq F_{S'}\in\IndCoh(S').
\]

\ob{2.3.6}{n:2.3.6}如同 §2.2.1，为了还原 $\Crys^r(Y)$，我们「不需要知道」DG 概形：

\begin{othm}{推论}{2.3.7}{n:2.3.7}
对 $Y\in\PreStk_{\mathrm{laft}}$，函子
\[
\Crys^r(Y)\to\varprojlim_{S\in((\Sch_{\mathrm{aff}}^{\mathrm{ft}})_{/Y_{\dR}})^{\mathrm{op}}}\IndCoh(S)
\]
是等价。
\end{othm}

\begin{proof}
容易由推论 1.3.7 推出。
\end{proof}

粗略地说，上述推论表明：$M\in\Crys^r(Y)$ 可以由如下资料还原：对每个 $S\in\Sch_{\mathrm{aff}}^{\mathrm{ft}}$ 及映射 $\clred S\to Y$，指定一个对象 $F_S\in\IndCoh(S)$；并且对每个 $f:S'\to S$，指定同构
\[
f^!(F_S)\simeq F_{S'}\in\IndCoh(S').
\]

\ob{2.3.8}{n:2.3.8}此外，推论 2.2.4 的类比对右晶体也成立：

\begin{othm}{推论}{2.3.9}{n:2.3.9}
设 $\mathcal{C}$ 是 (1.3) 中的任一范畴。则函子
\[
\Crys^r(Y)\to\varprojlim_{X\in(\mathcal{C}_{/Y})^{\mathrm{op}}}\Crys^r(X)
\]
是等价。
\end{othm}

\begin{proof}
由推论 1.3.9 得到。
\end{proof}

粗略地说，这个推论表明：对象 $M\in\Crys^r(Y)$ 可以由如下资料还原：对每个 $S\in\mathcal{C}_{/Y}$ 给定 $M_S\in\Crys^r(S)$，并对每个 $f:S'\to S$ 给定同构
\[
f^{\dagger,r}(M_S)\simeq M_{S'}\in\Crys^r(S').
\]

\ob{2.3.10}{n:2.3.10}自然变换 $p_{\dR}:\mathrm{Id}\to\dR$ 诱导自然变换
\[
\oblv^r:\Crys^r_{\PreStk_{\mathrm{laft}}}\to\IndCoh_{\PreStk_{\mathrm{laft}}}.
\]
也就是说，对每个 $Y\in\PreStk_{\mathrm{laft}}$ 有函子
\[
\oblv^r_Y:\Crys^r(Y)\to\IndCoh(Y),
\]
且对每个态射 $f:Y_1\to Y_2$ 有交换图
\begin{equation}
\begin{tikzcd}[column sep=large, row sep=large]
 \Crys^r(Y_1) \arrow[r, "\oblv^r_{Y_1}"] \arrow[d, "f^{\dagger,r}"'] & \IndCoh(Y_1) \arrow[d, "f^!"] \\
 \Crys^r(Y_2) \arrow[r, "\oblv^r_{Y_2}"'] & \IndCoh(Y_2).
\end{tikzcd}\tag{2.3}
\end{equation}
有：

\begin{othm}{引理}{2.3.11}{n:2.3.11}
若 $Y$ 经典形式光滑，则函子
\[
\Crys^r(Y)\to\Tot\big(\IndCoh(Y_\bullet/Y_{\dR})\big)
\]
是等价。
\end{othm}

\begin{proof}
与引理 2.1.7 的证明相同，即由引理 1.2.4 推出。
\end{proof}

\begin{othm}{引理}{2.3.12}{n:2.3.12}
对任意 $Y$，函子 $\oblv^r_Y$ 保守。
\end{othm}

\begin{proof}
由推论 2.3.9 与 (2.3) 的交换性，不妨设 $Y=X$ 是局部几乎有限型的仿射 DG 概形。取闭浸入 $i:X\to Z$，其中 $Z$ 是光滑经典有限型概形，并记 $Y$ 为 $Z$ 沿 $X$ 的形式完备化，$i':X\to Y$ 为相应的映射。

考虑交换图
\[
\begin{tikzcd}[column sep=large, row sep=large]
 \Crys^r(Y) \arrow[r, "\oblv^r_Y"] \arrow[d, "i'^{\dagger,r}"'] & \IndCoh(Y) \arrow[d, "i'^!"] \\
 \Crys^r(X) \arrow[r, "\oblv^r_X"'] & \IndCoh(X).
\end{tikzcd}
\]
此图中左竖直箭头是等价，因为 $i'_{\dR}:X_{\dR}\to Y_{\dR}$ 是同构。上方水平箭头保守，因为 $Y$ 形式光滑（特别地，经典形式光滑），可用引理 2.3.11。

于是只需说明函子 $i'^!$ 保守。这可以由 [GR1, Proposition 7.4.5] 与 [IndCoh, Proposition 4.1.7(a)] 合并得到。
\end{proof}

\section{左晶体与右晶体的比较}

请读者回想：我们假定所有预栈与 DG 概形都局部几乎有限型。

\ob{2.4.1}{n:2.4.1}回忆（见 [IndCoh, Sect.~5.7.5]）：对任意 $S\in\DGSch_{\mathrm{aff}}^{\mathrm{aft}}$，有典范定义的函子
\[
\Upsilon_S:\QCoh(S)\to\IndCoh(S),
\]
它由与对偶化层 $\omega_S\in\IndCoh(S)$ 作张量给出，使得对 $f:S_1\to S_2$，图表
\[
\begin{tikzcd}[column sep=large, row sep=large]
 \QCoh(S_1) \arrow[r, "\Upsilon_{S_1}"] \arrow[d, "f^*"'] & \IndCoh(S_1) \arrow[d, "f^!"] \\
 \QCoh(S_2) \arrow[r, "\Upsilon_{S_2}"'] & \IndCoh(S_2)
\end{tikzcd}
\]
典范交换。事实上，上述数据提升为函子的自然变换
\[
\Upsilon_{\DGSch_{\mathrm{aff}}^{\mathrm{aft}}}:\QCoh^*_{\DGSch_{\mathrm{aff}}^{\mathrm{aft}}}\to\IndCoh^!_{\DGSch_{\mathrm{aff}}^{\mathrm{aft}}},
\]
从而给出自然变换
\[
\Upsilon_{\PreStk_{\mathrm{laft}}}:\QCoh^*_{\PreStk_{\mathrm{laft}}}\to\IndCoh^!_{\PreStk_{\mathrm{laft}}},
\]
见 [IndCoh, Sect.~10.3.3]。

对单个对象 $Y\in\PreStk_{\mathrm{laft}}$，我们得到函子
\[
\Upsilon_Y:\QCoh(Y)\to\IndCoh(Y).
\]

\ob{2.4.2}{n:2.4.2}对 $Y\in\PreStk_{\mathrm{laft}}$，把 $\Upsilon$ 作用于 $Y_{\dR}$，得到典范定义的函子
\begin{equation}
\Upsilon_{Y_{\dR}}:\Crys^l(Y)\to\Crys^r(Y),\tag{2.4}
\end{equation}
使得图表
\begin{equation}
\begin{tikzcd}[column sep=large, row sep=large]
 \Crys^l(Y) \arrow[r, "\Upsilon_{Y_{\dR}}"] \arrow[d, "\oblv^l_Y"'] & \Crys^r(Y) \arrow[d, "\oblv^r_Y"] \\
 \QCoh(Y) \arrow[r, "\Upsilon_Y"'] & \IndCoh(Y)
\end{tikzcd}\tag{2.5}
\end{equation}
交换。

事实上，我们得到自然变换
\[
\Upsilon_{\PreStk_{\mathrm{laft}}}\circ\dR:\Crys^l_{\PreStk_{\mathrm{laft}}}\to\Crys^r_{\PreStk_{\mathrm{laft}}}.
\]
特别地，对 $f:Y_1\to Y_2$，图表
\[
\begin{tikzcd}[column sep=large, row sep=large]
 \Crys^l(Y_1) \arrow[r, "\Upsilon_{Y_{1,\dR}}"] \arrow[d, "f^{\dagger,l}"'] & \Crys^r(Y_1) \arrow[d, "f^{\dagger,r}"] \\
 \Crys^l(Y_2) \arrow[r, "\Upsilon_{Y_{2,\dR}}"'] & \Crys^r(Y_2)
\end{tikzcd}
\]
交换。

\ob{2.4.3}{n:2.4.3}我们断言：

\begin{othm}{命题}{2.4.4}{n:2.4.4}
对 $Y\in\PreStk_{\mathrm{laft}}$，函子 (2.4) 是等价。
\end{othm}

\begin{proof}
由推论 2.2.4 与 2.3.9，命题化归为：对几乎有限型的仿射 DG 概形 $X$，函子
\[
\Upsilon_{X_{\dR}}:\Crys^l(X)\to\Crys^r(X)
\]
是等价。

取闭浸入 $i:X\hookrightarrow Z$，其中 $Z$ 是光滑经典概形，并记 $Y$ 为 $Z$ 沿 $X$ 的形式完备化。由于 $X_{\dR}\to Y_{\dR}$ 是同构，函子
\[
f^{\dagger,l}:\Crys^l(Y)\to\Crys^l(X),\qquad f^{\dagger,r}:\Crys^r(Y)\to\Crys^r(X)
\]
都是等价。因此只需对 $Y$ 证明断言。

记 $Y_\bullet/Y_{\dR}$ 为 $\PreStk_{\mathrm{laft}}$ 中对应于映射
\[
p_{\dR,Y}:Y\to Y_{\dR}
\]
的 Čech 神经。

考虑交换图
\[
\begin{tikzcd}[column sep=large, row sep=large]
 \Crys^l(Y) \arrow[r, "\Upsilon_{Y_{\dR}}"] \arrow[d] & \Crys^r(Y) \arrow[d] \\
 \Tot\big(\QCoh(Y_\bullet/Y_{\dR})\big) \arrow[r, "\Tot(\Upsilon_{Y_\bullet/Y_{\dR}})"'] & \Tot\big(\IndCoh(Y_\bullet/Y_{\dR})\big).
\end{tikzcd}
\]
由引理 2.1.7 与 2.3.11，图中竖直箭头是等价。因此只需证明：对每个 $i$，
\[
\Upsilon_{Y^i/Y_{\dR}}:\QCoh(Y^i/Y_{\dR})\to\IndCoh(Y^i/Y_{\dR})
\]
是等价。

回忆（命题 1.3.3 的证明中也有）：$Y^i/Y_{\dR}$ 是光滑经典概形 $Z^i$ 沿 $X$ 的对角线副本的形式完备化。记 $U_i\subset Z^i$ 为补开子预栈。

由 [GR1, Propositions 7.1.3 and 7.4.5 and Diagram (7.16)]，我们得到 DG 范畴「短正合列」的映射
\[
\begin{tikzcd}[column sep=small, row sep=large]
 0 \arrow[r] & \QCoh(Y^i/Y_{\dR}) \arrow[r] \arrow[d, "\Upsilon_{Y^i/Y_{\dR}}"'] & \QCoh(Z^i) \arrow[r] \arrow[d, "\Upsilon_{Z^i}"] & \QCoh(U_i) \arrow[r] \arrow[d, "\Upsilon_{U_i}"] & 0 \\
 0 \arrow[r] & \IndCoh(Y^i/Y_{\dR}) \arrow[r] & \IndCoh(Z^i) \arrow[r] & \IndCoh(U_i) \arrow[r] & 0.
\end{tikzcd}
\]
现在，函子
\[
\Upsilon_{Z^i}:\QCoh(Z^i)\to\IndCoh(Z^i)\quad\text{与}\quad \Upsilon_{U_i}:\QCoh(U_i)\to\IndCoh(U_i)
\]
都是等价，因为 $Z^i$ 与 $U_i$ 光滑。

事实上，由 [IndCoh, Proposition 9.3.3]，对任意 $S\in\DGSch_{\mathrm{aff}}^{\mathrm{aft}}$，函子 $\Upsilon_S$ 是 $\Psi_S:\IndCoh(S)\to\QCoh(S)$ 的对偶；而后者在 $S$ 光滑时是等价，见 [IndCoh, Lemma 1.1.6]。
\end{proof}

\ob{2.4.5}{n:2.4.5}命题 2.4.4 使我们可以对 $Y\in\PreStk_{\mathrm{laft}}$ 认同左晶体与右晶体。

换言之，我们可以把范畴 $\Crys(Y)$ 看作带两个实现：「左」实现与「右」实现；它们分别体现为从 $\Crys(Y)$ 到 $\QCoh(Y)$ 与 $\IndCoh(Y)$ 的遗忘函子 $\oblv^l_Y$ 与 $\oblv^r_Y$。

两个遗忘函子由交换图
\begin{equation}
\begin{tikzcd}[column sep=large, row sep=large]
 & \Crys(Y) \arrow[dl, "\oblv^l_Y"'] \arrow[dr, "\oblv^r_Y"] & \\
 \QCoh(Y) \arrow[rr, "\Upsilon_Y"'] & & \IndCoh(Y)
\end{tikzcd}\tag{2.6}
\end{equation}
相联系。

对态射 $f:Y_1\to Y_2$，有自然定义的函子
\[
f^{\dagger}:\Crys(Y_2)\to\Crys(Y_1),
\]
它使下列图表交换：
\[
\begin{tikzcd}[column sep=large, row sep=large]
 \Crys(Y_1) & \Crys(Y_2) \arrow[l, "f^{\dagger}"'] \\
 \QCoh(Y_1) \arrow[u, "\oblv^l_{Y_1}"] & \QCoh(Y_2) \arrow[l, "f^*"] \arrow[u, "\oblv^l_{Y_2}"']
\end{tikzcd}
\]
以及
\[
\begin{tikzcd}[column sep=large, row sep=large]
 \Crys(Y_1) & \Crys(Y_2) \arrow[l, "f^{\dagger}"'] \\
 \IndCoh(Y_1) \arrow[u, "\oblv^r_{Y_1}"] & \IndCoh(Y_2) \arrow[l, "f^!"] \arrow[u, "\oblv^r_{Y_2}"'].
\end{tikzcd}
\]

\ob{2.4.6}{n:2.4.6}在下文中，我们将混用记号 $\Crys(Y)$、$\Crys^r(Y)$ 与 $\Crys^l(Y)$：前者强调命题不依赖于所选实现（左或右），后两者则在实现的选取重要时使用。

\ob{2.4.7}{n:2.4.7}引理 2.2.6 的证明。
\begin{proof}
由引理 2.3.12 与命题 2.4.4 合并得到。
\end{proof}

\section{Kashiwara 引理}

映射 $Y\mapsto\Crys(Y)$ 的一个特点是：Kashiwara 引理几乎成为同义反复。

我们将在晶体的右晶体实现中陈述并证明它。由命题 2.4.4，这也蕴含左晶体的相应断言。不过，用左晶体的语言写出同样的证明也很容易。

\ob{2.5.1}{n:2.5.1}回忆：$\PreStk$ 中的映射 $i:X\to Z$ 称为闭浸入，如果它在底层经典预栈的层面是闭浸入。也就是说，对每个 $S\in(\Sch_{\mathrm{aff}})_{/Z}$，基变换映射
\[
{}^{cl}(S\times_Z X)\to S
\]
是闭浸入；特别地，${}^{cl}(S\times_Z X)$ 是经典仿射概形。

若 $X,Z\in\PreStk_{\mathrm{laft}}$，则只需对 $S\in(\Sch_{\mathrm{aff}}^{\mathrm{ft}})_{/Z}$ 验证上述条件。

\ob{2.5.2}{n:2.5.2}对 $\PreStk_{\mathrm{laft}}$ 中对象的闭浸入 $i:X\hookrightarrow Z$，记 $j:Z^{\circ}\hookrightarrow Z$ 为补开浸入。诱导映射
\[
j_{\dR}:Z^{\circ}_{\dR}\to Z_{\dR}
\]
也是预栈的开浸入。考虑限制函子
\[
j^{\dagger,r}:\Crys^r(Z)\to\Crys^r(Z^{\circ}).
\]
由 [IndCoh, Lemma 4.1.1]，上述函子有全忠实的右伴随，记为 $j_{\dR,*}$，使得对每个 $S\in(\DGSch_{\mathrm{aff}}^{\mathrm{ft}})_{/Z_{\dR}}$ 及
\[
S^{\circ}:=S\times_{Z_{\dR}}Z^{\circ}_{\dR}\hookrightarrow S,
\]
由图
\[
\begin{tikzcd}[column sep=large, row sep=large]
 \IndCoh(S) & \Crys^r(Z) \arrow[l, "\oblv^r_Z"'] \\
 \IndCoh(S^{\circ}) \arrow[u, "(j_S)_*^{\IndCoh}"] & \Crys^r(Z^{\circ}) \arrow[l, "\oblv^r_{Z^{\circ}}"'] \arrow[u, "j_{\dR,*}"']
\end{tikzcd}
\]
经伴随性得到的图
\[
\begin{tikzcd}[column sep=large, row sep=large]
 \IndCoh(S) & \Crys^r(Z) \arrow[l, "\oblv^r_Z"'] \\
 \IndCoh(S^{\circ}) \arrow[u, "(j_S)^!"] & \Crys^r(Z^{\circ}) \arrow[l, "\oblv^r_{Z^{\circ}}"'] \arrow[u, "j^{\dagger,r}"']
\end{tikzcd}
\]
中的自然变换是同构。

特别地，自然变换
\[
\oblv^r_Z\circ j_{\dR,*}\to j_*^{\IndCoh}\circ\oblv^r_{Z^{\circ}}
\]
是同构。

\ob{2.5.3}{n:2.5.3}记 $\Crys^r(Z)_X$ 为 $\Crys^r(Z)$ 中等于 $\ker(j^{\dagger,r})$ 的全子范畴。

显然，对象 $M\in\Crys^r(Z)$ 属于 $\Crys^r(Z)_X$，当且仅当：对每个带映射 $\clred S\to Z$ 的 $S\in\DGSch_{\mathrm{aff}}^{\mathrm{ft}}$，相应的对象 $F_S\in\IndCoh(S)$ 落在
\[
\IndCoh(S)_{S-S^{\circ}}:=\ker\big(j_S^!:\IndCoh(S)\to\IndCoh(S^{\circ})\big)
\]
中。

\ob{2.5.4}{n:2.5.4}函子 $\Crys^r(Z)_X\hookrightarrow\Crys^r(Z)$ 有右伴随，由
\[
M\mapsto\Cone\big(M\to j_{\dR,*}\circ j^{\dagger,r}(M)\big)[-1]
\]
给出。因此，可以把 $\Crys^r(Z)_X$ 看作 $\Crys^r(Z)$ 的余局部化。

\ob{2.5.5}{n:2.5.5}由于复合 $i^{\dagger,r}\circ j_{\dR,*}$ 为零，函子 $i^{\dagger,r}:\Crys^r(Z)\to\Crys^r(X)$ 穿过上述余局部化：
\[
\Crys^r(Z)\to\Crys^r(Z)_X\xrightarrow{i'^{\dagger,r}}\Crys^r(X).
\]
Kashiwara 引理断言：

\begin{othm}{命题}{2.5.6}{n:2.5.6}
上述函子
\[
i'^{\dagger,r}:\Crys^r(Z)_X\to\Crys^r(X)
\]
是等价。
\end{othm}

\begin{proof}
注意在 $\PreStk_{\mathrm{laft}}$ 中有同构
\[
\varinjlim_{S\in(\DGSch_{\mathrm{aff}}^{\mathrm{ft}})_{/Z_{\dR}}}S\times_{Z_{\dR}}X_{\dR}\simeq X_{\dR}.
\]
此外，$S^{\wedge}:=S\times_{Z_{\dR}}X_{\dR}$ 等同于 $S$ 沿
\[
\clred S\times_{\clred Z}\clred X
\]
的形式完备化。

于是范畴 $\Crys^r(X)$ 可以描述为
\[
\varprojlim_{S\in((\DGSch_{\mathrm{aff}}^{\mathrm{ft}})_{/Z_{\dR}})^{\mathrm{op}}}\IndCoh(S^{\wedge}).
\]
按定义，范畴 $\Crys^r(Z)_X$ 由
\[
\varprojlim_{S\in((\DGSch_{\mathrm{aff}}^{\mathrm{ft}})_{/Z_{\dR}})^{\mathrm{op}}}
\ker\big(j_S^!:\IndCoh(S)\to\IndCoh(S^{\circ})\big)
\]
给出。

现在，[GR1, Proposition 7.4.5] 说：对上述任意 $S$，$!$-拉回给出等价
\[
\ker\big(j_S^!:\IndCoh(S)\to\IndCoh(S^{\circ})\big)\to\IndCoh(S^{\wedge}),
\]
这正是所要的。
\end{proof}

\begin{othm}{注}{2.5.7}{n:2.5.7}
若把上述证明改用左晶体而非右晶体的语言表述，我们就要用 [GR1, Proposition 7.1.3] 代替 [GR1, Proposition 7.4.5]。
\end{othm}
